\documentclass[a4paper,12pt]{article}

% Preamble einbinden
% --- Sprache und Kodierung ---
\usepackage[utf8]{inputenc}
\usepackage[T1]{fontenc}
\usepackage[ngerman]{babel}
\makeatletter
\newif\ifmodernlatex
\@ifl@t@r\fmtversion{2025-06-01}{\modernlatextrue}{\modernlatexfalse}
\makeatother

% --- Layout ---
\usepackage[a4paper,margin=2.5cm]{geometry}
\usepackage{setspace} % Zeilenabstand
\onehalfspacing
\usepackage{placeins} % Floatbarrier (Kapitel)

% --- Schriftarten ---
\usepackage{lmodern}

% --- Mathe und Symbole ---
\usepackage{amsmath, amssymb, amsthm}
\usepackage{calc}

% --- Grafiken ---
\usepackage{graphicx}
\ifmodernlatex
    \usepackage{tcolorbox}
    \tcbset{
        myboxstyle/.style={
            colframe=blue!60!black,
            colback=blue!5,
            boxrule=1pt,
            arc=2mm,
            left=3pt,right=3pt,top=3pt,bottom=3pt,
            width=\textwidth,
            fontupper=\footnotesize
        }
    }
\else
    \newenvironment{tcolorbox}[1][]
        {\par\medskip\begin{quote}\small}
        {\end{quote}\par\medskip}
\fi
\usepackage{float} % bessere Kontrolle von Abbildungen
\graphicspath{{figures/}}


% --- Tabellen ---
\usepackage{booktabs}
\usepackage{threeparttable}
\usepackage{tabularx}
\usepackage{makecell}
\usepackage{booktabs}
\ifmodernlatex
    \usepackage{siunitx}
    \sisetup{
        detect-all,
        table-format=3.2,
        separate-uncertainty = true,
        output-decimal-marker = {.},
        group-separator = {'}
    }
    \DeclareSIUnit{\gon}{gon}
\else
    \newcommand{\SI}[2]{\ensuremath{#1\,\mathrm{#2}}}
    \newcommand{\gon}{gon}
    \newcommand{\meter}{m}
\fi
\usepackage{multicol}
\usepackage{multirow}
\usepackage{xcolor}

% --- Referenzen und Links ---
\usepackage[hidelinks]{hyperref}
\usepackage{cleveref}
% Deutsche Bezeichnungen für \cref/\Cref
\crefname{figure}{Abbildung}{Abbildungen}
\Crefname{figure}{Abbildung}{Abbildungen}
\crefname{table}{Tabelle}{Tabellen}
\Crefname{table}{Tabelle}{Tabellen}
\crefname{equation}{Gleichung}{Gleichungen}
\Crefname{equation}{Gleichung}{Gleichungen}
\usepackage[labelfont=bf,font={footnotesize, it},labelsep=colon]{caption}
\usepackage[backend=biber,style=apa]{biblatex}
\addbibresource{bibliography.bib}
% Required for APA style
\DeclareLanguageMapping{english}{english-apa}


% --- Sonstiges ---
\usepackage{csquotes} % für saubere Zitate
\setlength{\marginparwidth}{2cm}
\usepackage[textsize=tiny]{todonotes}
\usepackage{ifthen}
\usepackage{xkeyval}
\usepackage{tikz}
\usepackage{lscape}
\usepackage{pdflscape}
\usepackage{amsmath}
\setcounter{tocdepth}{2}

\begin{document}

% Titelseite
\parindent 0mm
\begin{titlepage}
    \begin{figure}[t]
        \centering
        \includegraphics[width=16cm]{fhnw_habg_igeo_10mm.jpg}
    \end{figure}

    \vspace*{1.5cm}

    \begin{center}
        \textsf{
            \textbf{\LARGE Auswertung Netz Schwanderbärgli – V1.1: ‘optimiert’}\\[1cm]
            {\LARGE Teilgruppe 2 aus dem Feldkurs 2026}
        }
        \vspace{0.5cm}
    \end{center}

    \begin{figure}[h]
        \centering
        \includegraphics[width=16cm]{Titelbild.jpg}
    \end{figure}

    \vfill  % pushes the following text to the bottom

    \begin{flushleft}
        \textsf{
            \text{ }\\
            \textbf{Gruppe Brichen}\\[0.1cm]
            \textbf{Kelan Schmitter}\\
            Jascha Metzler\\
            Alina Portmann\\[0.5cm]
            Eingereicht bei: Dante Salvini - 5220 GeoSensorik \& Monitoring II\\
            \textbf{\today}\\
        }
    \end{flushleft}
\end{titlepage}


\newpage
\thispagestyle{empty}
\null
\vfill
%\pagenumbering{Roman}
\footnotesize{\textbf{Abbildung Titelseite:} Aufbau Netzmessung Gebiet Brichen} \\[2cm]
Version 1.0                                   \\[1cm]
Fachhochschule Nordwestschweiz (FHNW)         \\
Hochschule für Architektur, Bau und Geomatik  \\
Institut Geomatik                             \\
Hofackerstrasse 30                            \\
CH-4132 Muttenz                               \\
 
\verb"https://www.fhnw.ch"                   \\

\cleardoublepage

\input{chapters/00_zusammenfassung.tex}
\newpage

\tableofcontents
\newpage
\listoffigures
\listoftables
\newpage



% Kapitel laden
\input{chapters/01_Einleitung.tex}
\section{Literaturrecherche 3D-Netze}
\label{sec:literaturrecherche 3D-Netze}

\subsection{Einleitung}
Ein 3D-Netz bestimmt Lage und Höhe aller Netzpunkte in einer gemeinsamen Ausgleichung (Niemeier, 2008, S. 128, 247–248). 
GNSS liefert dabei 3D-Koordinaten oder Basislinien im geozentrischen System WGS84 (Niemeier, 2008, S. 330). 
Das Tachymeter (TPS) liefert Horizontalrichtungen, Zenitwinkel und Schrägdistanzen, die auf die lokale Lotrichtung bezogen sind (Banyai, 2011; França et al., 2021).
Die klassische Trennung in Lagenetz (2D) und Höhennetz (1D), wie sie etwa das Programm LTOP noch verwendet (swisstopo, o. J.), entfällt damit.
Der Bericht fasst zusammen, was die Literatur zu Aufbau, Modellierung, Ausgleichung und Qualitätsbeurteilung solcher Netze sagt. Hauptquellen sind das Lehrbuch von Niemeier (2008) zur Ausgleichungsrechnung und das Handbuch Ingenieurgeodäsie "Ingenieurbau" (Möser et al., [2016]).
Ergänzt wird mit Fachartikeln und Praxisbeispielen aus der Schweiz.
Niemeier (2008, S. 128) nennt zwei Wege, die Messungen mit dem Koordinatensystem in Einklang zu bringen: Reduktion auf ein Abbildungssystem (z. B. LV95, UTM) oder Formulierung der Beobachtungsgleichungen direkt in einem räumlichen System. 
Für hochgenaue lokale Ingenieurnetze ist ein 3D-System von Bedeutung, wenn die Höhe streng mitverarbeitet werden soll. Für die Kombination von GNSS mit terrestrischen Messungen ist das globale kartesische System gebräuchlich.

\subsection{Messgrössen des 3D-Netzes}
\label{sec:Messgrössen des 3D-Netzes}
GNSS und TPS ergänzen sich folgendermassen: GNSS bestimmt Lagerung, Orientierung und Massstab des Netzes, Strecken- und Richtungsmessungen liefern die innere Geometrie mit hoher relativer Genauigkeit (Niemeier, 2008, S. 232, 247–248). 
Im Tunnelbau wird zum Beispiel ein GNSS-gestütztes Portalnetz mit Tachymetermessungen in den Berg übertragen (Leica Geosystems, 2016; Weber, 2014). 
Die folgende Tabelle zeigt die Messgrössen mit typischen Herstellerangaben aktueller Präzisionsinstrumente.

\begin{table}[h!]
    \centering
    \label{tab:messgrössen}
    \begin{tabularx}{\textwidth}{*{3}{>{\centering\arraybackslash}X}}
        \toprule
        \textbf{Messgrössen} & \textbf{Bezugssystem} & \textbf{Herstellergenauigkeit} \\
        \midrule
        GNSS-Basislinie $\Delta X$, $\Delta Y$, $\Delta Z$ (statisch) & geozentrisch (WGS84/ETRS89) & Lage 3 mm + 0,5 ppm, Höhe 5 mm + 0,5 ppm (Leica GS18) \\
        Horizontalrichtung (TPS) & lokal astronomisch (Lotrichtung) & $0{,}5^{\prime\prime} \approx 0{,}15\,\mathrm{mgon}$ (Leica TS60) \\
        Zenitwinkel (TPS) & lokal astronomisch (Lotrichtung) & $0{,}5^{\prime\prime} \approx 0{,}15\,\mathrm{mgon}$ (Leica TS60) \\
        Schrägdistanz (TPS) & systemunabhängig & 0,6 mm + 1 ppm (Leica TS60) \\
        Höhenunterschied (Nivellement) & Schwerefeld (Gebrauchshöhen) & Feldkurs Trimble Niv. NOCH Nachschauen \\
        \bottomrule
    \end{tabularx}
    \caption{Beispielgrössen in 3D-Netzen.}
\end{table}

Niemeier (2008, S. 330) betont, dass GNSS-Auswerteprogramme 3D-Koordinaten oder Basislinien im WGS84 liefern, die meist schon als Näherungswerte für die Netzausgleichung taugen.
Wichtig ist, die von der Software gelieferten voll besetzten Kofaktormatrizen in das stochastische Modell zu übernehmen, statt sie durch Diagonalmatrizen zu ersetzen (Niemeier, 2008, S. 249).
Die geringere Genauigkeit der GNSS-Höhenkomponente lässt sich so streng berücksichtigen (Niemeier, 2008, S. 365).
Die Herstellerangaben gelten für Idealbedingungen.
Im Netz kommen Zentrierfehler, Refraktion und Lotabweichung hinzu, die im funktionalen und stochastischen Modell berücksichtigt werden müssen (vgl. Abschnitt 3).

\subsection{Bezugssysteme und Modellbildung}
\label{sec:Bezugssysteme und Modellbildung}

\subsection{Ausgleichung im Gauss-Markov-Modell}
\label{sec:Ausgleichung im Gauss-Markov-Modell}

\subsection{Genauigkeit und Zuverlässigkeit}
\label{sec:Genauigkeit und Zuverlässigkeit}

\subsection{Fazit?}
\label{sec:Fazit?}
\clearpage
\section{Datenaufbereitung TPS}
\label{sec:tps}
Dieses Kapitel beschreibt, wie die im Feldkurs erhobenen tachymetrischen Messungen
für die Netzausgleichung aufbereitet wurden.

\subsection{Messkonfiguration und Instrumente}
\label{sec:messkionfiguration und instrumente}
Die tachymetrischen Messungen wurdem mit drei Leica Nova MS60 (interne Nr. 5, 6 und 8)
und einem Leica Nova TS60 (interne Nr. 2) durchgeführt 

\subsection{Datenübertragung und -sicherung}
\label{sec:datenuebertragung und -sicherung}
Die Messdaten der gesamten Feldkampagne wurden redundant auf MS Teams und auf einer
externen Festplatte (G5 Nr. 4) gesichert. Die Rohdaten wurden unverändert und getrennt
von den bearbeiteten Daten abgelegt.

\subsection{Vorgehen in Leica Infinity}
\label{sec:vorgehen in leica infinity}
Die Aufbereitung erfolgte in Leica Infinity. Pro Station wurden folgende Schritte durchgeführt:
\begin{enumerate}
    \item Import der Messdaten
    \item Kontrolle der gemessenen Visuren auf Vollständigkeit und Plausibilität
    \item Prüfung der Satzmittel und Standardabweichungen
    \item Deaktivierung einzelner Sätze bei Auffälligkeiten (Kapitelxxx)
    \item Kontrolle und Anpassung der Meteokorrekturen
    \item Nachbearbeitung einzelner Stationen
\end{enumerate}

\subsection{Metekorrekturen}
\label{sec:metekorrekturen}
Bei der Stationierung wurden die Meteowerte am Standpunkt eingegeben. Das Netz
überspannt jedoch einen Höhenunterschied von rund 570 m, sodass diese Werte die
Verhältnisse entlang langer Visuren nur bedingt abbilden. Die Meteokorrekturen wurden
daher bei langen Visuren überprüft und angepasst. In Leica Infinity können dafür auch
die Meteowerte der Zielpunkte berücksichtigt werden. Beim Standpunkt 6003 wurden die
Meteokorrekturen zusätzlich angebracht

\subsection{Stationsspezifische Bearbeitung}
\label{sec:stationsspezifische_bearbeitung}
Die folgenden Abschnitte listen alle Stationen auf, bei denen Sätze deaktiviert oder
Daten angepasst wurden.

\subsubsection{Gebiet Sperre}
\label{sec:gebiet_sperre}
Das Teilgebiet Sperre umfasst die obere Schwergewichtsmauer. Alle Profilpunkte
der Mauer konnten vom Standpunkt 1000 aus gemessen werden. Die Punkte 340, 520, 530, 540, 720, 730,
740 wurden nur von einer Station eingemessen.

\begin{table}[H]
    \centering
    \begin{tabular}{p{1.5cm} p{6cm} p{5.5cm}}
        \toprule
        Station & Feststellung & Massnahme/Bemerkung \\
        \midrule
        1000 & Punkt 210 wurde nur im 1. Satz gemessen. Im 2. Satz teils
        grössere Abweichungen (noch akzeptabel). Einzelne Punkte sind in
        Infinity nicht deaktivierbar, deshalb wurde Punkt 210 vom 210 ausgeschlossen.
        Alle Sätze wurden verwendet.
        & Nachmessung auf 210 in .mes-Datei eingefügt (Erklärung siehe unten) \\
        \addlinespace
        2000 & In der vorhandenen .mes-Datei fehlte die Beobachtung nach
        701. Sie war bei der Datenerfassung vergessen gegangen.
        & Der gesamte Standpunkt 2000 wurde aus den Daten des .mes-Files der SX12 ergänzt. \\
        \bottomrule
    \end{tabular}
    \caption{Feststellungen und Massnahmen bei der Datenaufbereitung Gebiet Sperre}
    \label{tab:feststellungen}
\end{table}

\textbf{Station 1000:} Die Nachmessung der Visur $1000 \rightarrow 210$ wurde im Feld
nur mit der Anwendung «Messen» und nicht in der Satzmessung aufgenommen. In der
Auswertung gehen jedoch Richtungen im System der Satzmessungen ein. Die Messung wurde
daher zuerst umgerechnet und anschliessend als Richtung RI210 in die .mes-Datei eingefügt.

Dazu wurde die Orientierung der Satzmessung zum Punkt 710 bestimmt, der sowohl in der
Satzmessung als auch in «Messen» gemessen wurde. Sie ergibt sich als Differenz zwischen
dem Azimut und der reduzierten Richtung der Satzmessung:
\begin{align*}
    Orientierung_{710} &= \SI{338.0570}{\gon} - \SI{348.5455}{\gon} = \SI{-10.4885}{\gon}
\end{align*}

Zur Kontrolle wurde die Orientierung zusätzlich zu Punkt 610 berechnet:
\begin{align*}
    Orientierung_{610} &= \SI{347.4518}{\gon} - \SI{357.9403}{\gon} = \SI{-10.4885}{\gon}
\end{align*}

Beide Punkte liefern denselben Wert. Die in «Messen» gemessene Richtung auf Punkt 210
($r_{\text{Messen}} = \SI{396.45123}{\gon}$) wurde anschliessend um die Orientierung
korrigiert und auf den Wertebereich von $0$ bis $\SI{400}{\gon}$ reduziert:
\begin{align*}
    r_{210} &= \SI{396.45123}{\gon} - (\SI{-10.4885}{\gon}) = \SI{406.93973}{\gon} \\
    r_{210} &= \SI{406.93973}{\gon} - \SI{400}{\gon} = \SI{6.93973}{\gon}
\end{align*}

\subsubsection{Gebiet Brichen}
\label{sec:gebiet-brichen}
Das Teilgebiet Brichen konnte im Feld gemäss Messkonzept ohne
Auffälligkeiten gemessen werden (Wälti et al., 2026).
Der im Feld fehlende Punkt 7019 wurde neu materialisiert und unter der
Nummer 7119 eingemessen. Bei der Datenaufbereitung ergaben sich die
folgenden Feststellungen.

\begin{table}[H]
    \centering
    \begin{tabularx}{\textwidth}{p{2.2cm} X X}
        \toprule
        Station & Feststellung & Massnahme/Bemerkung \\
        \midrule
        7015\_1, 7015\_2 & Von Station 7015\_1 wurden zwei Punkte
        angemessen (7014, EX7016.4), von Station 7015\_2 nur ein Punkt
        (7014). & Beide Stationen wurden verwendet. \\
        \addlinespace
        7016\_2 & Die Sätze 1--3 hatten grössere Abweichungen.
        & Es wurden nur die Sätze 4--6 verwendet. \\
        \addlinespace
        7119 & Der erste Satz hatte grössere Abweichungen. &
        Es wurden nur die Sätze 2--6 verwendet. \\
        \addlinespace
        7023 & Punkt EX7016.4 wurde mit der falschen Punktnummer
        (EX7016.2) gemessen. & Im .mes-File wurde die Punktnummer
        korrigiert. \\
        \bottomrule
    \end{tabularx}
    \caption{Feststellungen und Massnahmen bei der Datenaufbereitung
    Gebiet Brichen}
    \label{tab:feststellungen-brichen}
\end{table}

\subsubsection{Gebiet Ägerdi}
\label{sec:gebiet-aegerdi}
Das Gebiet Ägerdi bildet den obersten Teil des Überwachungsperimeters.
Die Punkte liegen auf einem Polygonzug entlang der Strasse und verteilt
über die angrenzende Wiese. Zusätzlich wird ein Gebäude über eine
Zielmarke (6050) überwacht. Die Messungen erfolgten am zweiten Tag der
Messkampagne und konnten nach der Rodung einzelner Visuren wie geplant
ausgeführt werden, ergänzt um eine zusätzliche Visur. Der durch eine
Strassensanierung zerstörte Punkt 6019 wurde durch einen neuen Punkt
ersetzt und unter der Nummer 6119 ins Netz aufgenommen (Wälti et al., 2026). Bei der Datenaufbereitung
ergaben sich die folgenden Feststellungen.

\begin{table}[htbp]
    \centering
    \begin{tabularx}{\textwidth}{p{2.2cm} X X}
        \toprule
        Station & Feststellung & Massnahme/Bemerkung \\
        \midrule
        6001 & Die Standardabweichungen der Satzmittel sind etwas höher
        als bei den übrigen Stationen. & Die Werte liegen innerhalb der
        Gerätegenauigkeit. Es wurden keine Anpassungen vorgenommen. \\
        \addlinespace
        6003 & Die Sätze 1--2 wiesen erhöhte Abweichungen auf. &
        Es wurden nur die Sätze 3--6 verwendet.
        Zusätzlich wurden die Meteokorrekturen angebracht. \\
        \addlinespace
        6016 & Der erste Satz wies grössere Abweichungen auf &
        Es wurden nur die Sätze 2--6 verwendet. \\
        \addlinespace
        7001 & Der erste Satz wies grössere Abweichungen auf. &
        Es wurden nur die Sätze 2--6 verwendet. \\
        \addlinespace
        TP506 & Der dritte Satz wies grössere Abweichungen auf. &
        Es wurden die Sätze 1--2 und 4--6 verwendet. \\
        \bottomrule
    \end{tabularx}
    \caption{Feststellungen und Massnahmen bei der Datenaufbereitung
    im Gebiet Ägerdi}
    \label{tab:feststellungen-aegerdi}
\end{table}

\subsubsection{Gebiet Brunni}
\label{sec:gebiet-brunni}
Brunni ist ein kleines Teilgebiet auf einer Waldlichtung. Die steile
Wiese ist nur zu Fuss zugänglich, und alle Überwachungspunkte liegen auf
der Wiese. Die Verbindung zu den übrigen Teilgebieten ist nur über GNSS
oder über einzelne Visuren zu den Punkten 6000 und 7022 möglich (Wälti et al., 2026). Das Gebiet wurde am zweiten Tag der
Messkampagne zusammen mit dem Gebiet Ägerdi gemessen und konnte wie
geplant vermessen werden. Zusätzlich wurde vom Punkt 7003 eine Visur
zum Punkt 7022 und vom Punkt 6011 eine Visur zum Punkt 6000 gemessen
(Wälti et al., 2026). Bei der Datenaufbereitung waren im
Gebiet Brunni keine Anpassungen an den Messdaten nötig.

\subsection{Ergebnis der Datenaufbereitung}
\label{sec:ergebnis_der_datenaufbereitung}

Nach Abschluss der beschriebenen Kontrollen und Anpassungen liegt ein
bereinigter tachymetrischer Datensatz für die weitere Netzausgleichung
vor. Auffällige Messsätze wurden stationsweise beurteilt und, wo
erforderlich, von der weiteren Auswertung ausgeschlossen. Fehlende oder
fehlerhaft bezeichnete Beobachtungen wurden anhand der vorhandenen
Rohdaten ergänzt beziehungsweise korrigiert. Zudem wurden die
Meteokorrekturen insbesondere bei langen Visuren überprüft und
angepasst. Das resultierende.mes-File bildet damit das
Ergebnis der TPS-Datenaufbereitung und die Grundlage für die
nachfolgenden Netzausgleichungen.

Zur Plausibilitätskontrolle des aufbereiteten Datensatzes wurde eine
erste weich gelagerte Netzausgleichung mit LTOP durchgeführt. Diese
diente nicht der abschliessenden Auswertung des Netzes, sondern der
Kontrolle, ob nach der Datenaufbereitung noch grobe Fehler,
Punktverwechslungen oder auffällige Beobachtungen vorhanden sind.
Die Resultate zeigen, dass der TPS-Datensatz grundsätzlich konsistent
weiterverwendet werden kann.

Eine noch zu untersuchende Auffälligkeit zeigt sich in der
Höhenkomponente der langen Visur zwischen 510 und EX9500. Bei der rund
1.26~km langen Verbindung ergibt sich eine Abweichung von rund 52~mm.
Da die Visur einen grossen Höhenunterschied überwindet, sind bei der
weiteren Auswertung insbesondere mögliche atmosphärische Einflüsse auf
die Höhenübertragung zu untersuchen. Eine abschliessende Beurteilung
dieser Auffälligkeit ist allein anhand der im vorliegenden Arbeitspaket
aufbereiteten TPS-Daten nicht möglich.

Im Feldkurs wurden für ausgewählte Punkte zusätzlich langstatische
GNSS-Messungen, eine trigonometrische Höhenbestimmung sowie ein
Präzisionsnivellement durchgeführt. Diese unabhängigen Messungen werden
in separaten Arbeitspaketen ausgewertet und können später zur Kontrolle
und Kombination der unterschiedlichen Höheninformationen herangezogen
werden. Sie sind daher nicht Bestandteil der vorliegenden
TPS-Datenaufbereitung.

Mit dem bereinigten .mes-File ist die TPS-Datenaufbereitung
abgeschlossen. Die detaillierte Beurteilung der Beobachtungen sowie
deren Kombination mit den übrigen Messverfahren erfolgt in den
nachfolgenden Arbeitspaketen zur Netzausgleichung.
\section{Auswertung Netz Schwanderbärgli optimiert V1.1}
\label{sec:auswertung netz schwanderbärgli optimiert v1.1}
Im Anschluss an die Datenaufbereitung wurden die Messdaten der
Feldkampagne 2026 mit dem Programm GeoSuite (swisstopo LTOP)
ausgewertet. Die bereinigten Beobachtungen wurden im `.mes`-Dateiformat
und die Näherungskoordinaten der Netzpunkte im `.koo`-Dateiformat in
LTOP eingelesen. Als Grundlage für die Näherungskoordinaten dienten
die Koordinaten aus dem Jahr 2022. Für die im Rahmen der Feldkampagne
2026 mittels GNSS bestimmten Punkte wurden die aktuellen
GNSS-Koordinaten übernommen. Die Netzausgleichung erfolgte nach der
Methode der kleinsten Quadrate (MdKQ). Dabei wird angenommen, dass die
Beobachtungen stochastisch unabhängig und die Messabweichungen
normalverteilt sind.

Im Rahmen dieses Arbeitspakets wird das Netz Schwanderbärgli in der
Variante V1.1 «optimiert» ausgewertet. Ziel ist es, durch eine geeignete
Auswahl und Kombination der vorhandenen tachymetrischen Messungen (TPS),
GNSS-Beobachtungen, Nivellementdaten und trigonometrischen
Höhenbestimmungen eine möglichst einfache Netzkonfiguration zu
entwickeln, ohne die erforderliche Genauigkeit und Zuverlässigkeit
wesentlich zu beeinträchtigen. Dazu werden verschiedene
Netzkonfigurationen und Lagerungsarten untersucht und hinsichtlich
ihrer Koordinatengenauigkeiten, Beobachtungsverbesserungen und
Zuverlässigkeitskennwerte miteinander verglichen. Anhand dieser
Ergebnisse soll beurteilt werden, welche Beobachtungen für eine
zukünftige, vereinfachte Überwachung des Schwanderbärgli weiterhin
erforderlich sind.

\subsection{Provisorischer Abriss TPS, GNSS und Nivellement}
\label{sec:provisorischer abriss tps gnss und nivellement}
Als erster Schritt wurde ein provisorischer Abriss erstellt, um eine erste
Qualitätskontrolle des Netzes durchzuführen. Beim provisorischen Abriss werden
alle Punkte als Festpunkte behandelt. Es handelt sich somit um eine reine
Plausibilitätskontrolle der gemessenen Richtungen und Distanzen anhand der
bekannten Näherungskoordinaten.

Der provisorische Abriss zeigte insgesamt ein plausibles und stimmiges Messnetz.
Die Hin- und Rückmessungen der Distanzen stimmen sehr gut überein (grösste
Differenz 5 mm bei 41 Distanzpaaren), und auch die gegenseitig gemessenen
Höhenwinkel zeigen im Median keine Abweichungen. Es wurden keine groben Fehler
in den Beobachtungen festgestellt. Einzelne Beobachtungen zeigten grössere
Abweichungen. Diese treten gruppiert auf, etwa als gemeinsame Verschiebung der
Distanzen ab Station 2000, und sind durch Bewegungen im Messgebiet sowie durch
die verschiedenen Näherungskoordinaten (GNSS Koordinaten 2026 und Koordinaten 2022) erklärbar.
Sie blieben somit im erwarteten Bereich eines provisorischen Abrisses. Die
Auswertung ergab 360 Lagebeobachtungen (davon 166 Richtungen, 166 Distanzen und
28 GNSS-Koordinaten) sowie 183 Höhenbeobachtungen (davon 169 Höhendifferenzen
und 14 GNSS-Höhen) bei 118 Stationen.

Eine Auffälligkeit zeigte die gegenseitige Höhenkontrolle zwischen den Punkten
510 und EX9500. Dort beträgt die Differenz der Hin- und Rückmessung 53.5 mm bei
einer Distanz von rund 1264 m. Da zum aktuellen Zeitpunkt noch trigonometrische
Höhenbestimmungen fehlen, wird diese Beobachtung vorerst nicht weiter beurteilt.
Das Problem wird zu einem späteren Zeitpunkt angeschaut.

\subsection{Weiche Lagerung TPS, GNSS und Nivellement}
\label{sec:weiche lagerung tps gnss und nivellement}
Nach dem provisorischen Abriss wurde das Netz in einer weichen Lagerung
ausgeglichen. Auf eine vorgängige freie Netzausgleichung wurde verzichtet.
Stattdessen wurden die vorhandenen Referenzkoordinaten bereits in der
ersten Ausgleichung berücksichtigt. Durch die weiche Lagerung können
diese innerhalb ihrer vorgegebenen Genauigkeiten angepasst werden,
ohne dass das Netz durch eine starre Festhaltung der Referenzpunkte
unnötig gezwängt wird. Dies ist insbesondere bei einem ausgedehnten
Überwachungsnetz mit unterschiedlichen Beobachtungstypen und
Referenzkoordinaten aus verschiedenen Messepochen von Bedeutung.

Ziel der weichen Lagerung war es, die innere Konsistenz der
Beobachtungen zu überprüfen, auffällige Messungen zu identifizieren
und eine geeignete Ausgangslage für die weitere Netzoptimierung
zu schaffen. Dazu wurden insbesondere die Verbesserungen der
Beobachtungen sowie die Genauigkeits- und Zuverlässigkeitskennwerte
der Ausgleichung untersucht.
Die Optimierung erfolgte schrittweise in den Varianten
weiche Lagerung v0 bis v5. Dabei wurden auffällige Beobachtungen nicht
allein aufgrund der statistischen Kennwerte eliminiert, sondern
zusätzlich anhand der Messgeometrie und der bekannten Bedingungen
während der Feldmessung beurteilt.

\subsubsection{V0}
\label{sec:v0}

\section{Scanning}
\label{sec:scanning}

\subsection{Datenaufbereitung}
\label{sec:datenaufbereitung}

\subsection{Datenauswertung}
\label{sec:datenauswertung}

\subsection{Resultate}
\label{sec:resultate}

\subsection{Diskussion der Resultate}
\label{sec:diskussion der resultate}
\section{Scanning}
\label{sec:scanning}

\subsection{Datenaufbereitung}
\label{sec:datenaufbereitung}

\subsection{Datenauswertung}
\label{sec:datenauswertung}

\subsection{Resultate}
\label{sec:resultate}

\subsection{Diskussion der Resultate}
\label{sec:diskussion der resultate}
\section{Reflexion}
\label{sec:reflexion}

Dies ist Platzhaltertext für das Reflexion. Hier werden die wichtigsten Erkenntnisse des
Berichts kurz zusammengefasst und ein Ausblick auf allfällige weitere Schritte gegeben.

\section{Beispiele}

Dies ist Platzhaltertext für die Diskussion der Resultate. Hier werden die in
Kapitel~\ref{sec:resultate} vorgestellten Ergebnisse eingeordnet und kritisch bewertet.

\subsection{Beispiel für Querverweise}
\label{subsec:querverweise}
Mit dem Paket \texttt{cleveref} können Querverweise automatisch mit dem passenden Typ
beschriftet werden, zum Beispiel ein Verweis auf \Cref{tab:beispieltabelle} oder auf
\Cref{fig:beispielbild}. Ebenso kann erneut auf eine Quelle verwiesen werden
\cite{muster_beispielquelle_2025}.

\subsection{Einschränkungen}
\label{subsec:einschraenkungen}
Dies ist Platzhaltertext für die Beschreibung von Einschränkungen oder offenen Fragen.


\subsection{Beispiel mit farbig hervorgehobener Tabelle}
\label{subsec:beispiel-farbtabelle}
Mit dem Paket \texttt{xcolor} (Option \texttt{table}) können einzelne Tabellenzeilen
hervorgehoben werden, siehe Tabelle~\ref{tab:beispiel-farbtabelle}.

\begin{table}[h!]
    \centering
    \caption{Beispieltabelle mit farbig hervorgehobenen Zeilen.}
    \label{tab:beispiel-farbtabelle}
    \begin{tabularx}{\textwidth}{l *{2}{>{\centering\arraybackslash}X}}
        \toprule
        \textbf{Punktnr.} & \textbf{Abweichung [cm]} & \textbf{Bemerkung} \\
        \midrule
        1001 & 1.2 & Platzhalter \\
        1002 & 3.4 & \textcolor{red}{\textbf{Auffälliger Wert}} \\
        1003 & 0.8 & Platzhalter \\
        \bottomrule
    \end{tabularx}
\end{table}

\subsection{Beispiel einer mehrspaltigen Liste}
\label{subsec:beispiel-multicol}
Mit dem Paket \texttt{multicol} können Aufzählungen auf mehrere Spalten verteilt werden,
zum Beispiel für Listen mit vielen kurzen Einträgen: \todo{Dies ist ein To-do Kommentar}

\begin{multicols}{2}
    \begin{itemize}
        \item Eintrag 1
        \item Eintrag 2
        \item Eintrag 3
        \item Eintrag 4
        \item Eintrag 5
        \item Eintrag 6
    \end{itemize}
\end{multicols}

\subsection{Beispiel eines einfachen Diagramms}
\label{subsec:beispiel-tikz}
Mit dem Paket \texttt{tikz} können einfache Diagramme direkt in \LaTeX{} erstellt werden,
siehe \Cref{fig:beispiel-tikz}.

\begin{figure}[h!]
    \centering
    \begin{tikzpicture}
        \node[draw, rectangle, minimum width=2.5cm, minimum height=1cm] (a) at (0,0) {Schritt 1};
        \node[draw, rectangle, minimum width=2.5cm, minimum height=1cm] (b) at (4,0) {Schritt 2};
        \draw[->] (a) -- (b);
    \end{tikzpicture}
    \caption{Beispiel für ein einfaches \texttt{tikz}-Diagramm.}
    \label{fig:beispiel-tikz}
\end{figure}

\FloatBarrier


\subsection{Zielsetzung}
Dies ist eine Aufzählung:
\begin{itemize}
    \item 1 Punkt.
    \item 2 Punkt.
    \item ...
\end{itemize}


\subsection{Vorgehen}
\label{subsec:vorgehen}
Dies ist Platzhaltertext zur Beschreibung des methodischen Vorgehens. Hier werden die
einzelnen Arbeitsschritte erläutert.

\begin{tcolorbox}[myboxstyle, title=Hinweis]
Dies ist eine Beispiel-Hinweisbox. Sie eignet sich, um wichtige
Hinweise, Annahmen oder Einschränkungen hervorzuheben.
\end{tcolorbox}

\subsection{Beispielrechnung}
\label{subsec:beispielrechnung}
Formeln können mit \texttt{amsmath} gesetzt werden, zum Beispiel die Berechnung einer
Bodenauflösung (GSD) in Abhängigkeit der Flughöhe:
\begin{equation}
    \text{GSD} = \frac{h \cdot p}{f}
    \label{eq:gsd}
\end{equation}
mit Flughöhe $h$, Pixelgrösse $p$ und Kamerakonstante $f$. Mit dem Befehl \texttt{\textbackslash SI}
können Einheiten konsistent gesetzt werden, zum Beispiel eine Flughöhe von \SI{60}{\meter}.
Gleichung~\ref{eq:gsd} bzw.\ \cref{eq:gsd} kann so im Text referenziert werden.

\subsection{Beispieltabelle}
\label{subsec:beispieltabelle}
Tabelle~\ref{tab:beispieltabelle} zeigt den Aufbau einer Beispieltabelle mit
\texttt{tabularx}, \texttt{threeparttable}, \texttt{multirow} und \texttt{makecell}.

\begin{table}[h!]
    \centering
    \caption{Kontrollmessungen zur Bestimmung des Fehlers aus den Basisdaten.}
    \label{tab:beispieltabelle}
    \begin{tabularx}{\textwidth}{c | *{3}{>{\centering\arraybackslash}X}}
        \toprule
        \makecell{\textbf{Höhe} \\ \textbf{[m.ü.M]}} & 
        \makecell{$\mathbf{\Delta E_{VRS} - E_{MDR}}$ \\ \textbf{[cm]}} & 
        \makecell{$\mathbf{\Delta N_{VRS} - N_{MDR}}$ \\ \textbf{[cm]}} & 
        \makecell{$\mathbf{\Delta H_{VRS} - H_{MDR}}$ \\ \textbf{[cm]}} \\
        \midrule
        1885 (Fix 1) & -2.1 & 0.8 & \textcolor{red}{\textbf{5.5}} \\
        1885 (Fix 2) & -2.6 & 1.8 & \textcolor{red}{\textbf{5.2}} \\
        2085 & -0.6 & -2.5 & \textcolor{red}{\textbf{3.6}} \\
        2145 & -0.7 & -2.9 & \textcolor{red}{\textbf{5.9}} \\
        2205 & -0.5 & -2.4 & \textcolor{red}{\textbf{5.1}} \\
        2255 & -1.1 & -2.1 & \textcolor{red}{\textbf{4.6}} \\
        \midrule
        \textbf{} & 
        $\mathbf{RMSE_E}$ \textbf{[cm]} & 
        $\mathbf{RMSE_N}$ \textbf{[cm]} & 
        $\mathbf{RMSE_H}$ \textbf{[cm]} \\
        \midrule
        \textbf{} & 0.9 & 2.4 & \textcolor{red}{\textbf{5.0}} \\
        \bottomrule
    \end{tabularx}
\end{table}

\subsection{Beispielabbildung}
\label{subsec:beispielabbildung}
\Cref{fig:beispielbild} zeigt eine Beispielabbildung inklusive Bildlegende und Quellenangabe.

\begin{figure}[h!]
    \centering
    \includegraphics[width=\textwidth]{aufnahmeperimeter_Skipiste1.jpg}
    \caption[Beispielabbildung zur Illustration des Bild- und Beschriftungsformats]{Beispielabbildung zur Illustration des Bild- und Beschriftungsformats.
    \\Projektion:~LV95~|~Stand:~TT.MM.JJJJ~|~Datenquelle:~Platzhalter}
    \label{fig:beispielbild}
\end{figure}

\FloatBarrier




% Literatur
\newpage
\section{Quellenverzeichnis}
\printbibliography

\section{Anhang}
\label{sec:anhang}

Dies ist Platzhaltertext für den Anhang, zum Beispiel eine Liste der beigelegten
Dokumente:
\begin{itemize}
    \item 01\_Beispieldokument
    \item 02\_Beispieldokument
\end{itemize}

\subsection{Beispiel einer verschachtelten Datenstruktur}
\label{subsec:datenstruktur}
Verschachtelte Aufzählungen eignen sich zum Beispiel zur Darstellung einer
Ordnerstruktur:
\begin{itemize}
    \item Projektordner
    \begin{itemize}
        \item 10\_Rohdaten
        \begin{itemize}
            \item Beispieldatei\_1
            \item Beispieldatei\_2
        \end{itemize}
        \item 20\_Auswertung
        \item 30\_Ausgabe
    \end{itemize}
\end{itemize}


\end{document}
